\section{第六章：变分法}

\begin{frame}{第六章：从求一个数到求一条函数}
  普通极值问题寻找使 $f(x)$ 极值的数 $x$；变分问题寻找使泛函 $J[u]$ 极值或驻定的函数 $u(x)$。
  \[
    f(x)\quad\longrightarrow\quad J[u].
  \]
  \begin{block}{典型泛函}
    \[
      J[u]=\int_a^bL(x,u,u')\,\mathrm{d}x.
    \]
  \end{block}
  它可表示曲线长度、运动时间、作用量或系统能量。
\end{frame}

\begin{frame}{例：两点之间的最短曲线}
  曲线 $y=u(x)$ 的弧长是
  \[
    J[u]=\int_{x_1}^{x_2}\sqrt{1+u'(x)^2}\,\mathrm{d}x,
    \qquad u(x_1)=y_1,\quad u(x_2)=y_2.
  \]
  \begin{itemize}
    \item 未知对象不是一个点，而是满足端点条件的所有曲线。
    \item 满足条件的曲线称为\blue{可容许函数}。
    \item 端点条件是问题的一部分；端点不同的曲线不能相互比较。
  \end{itemize}
  \begin{alertblock}{目标}
    找到使长度的一阶变化为零的路径；结果将恢复熟悉的直线。
  \end{alertblock}
\end{frame}

\begin{frame}{什么是变分？}
  从试探函数 $u(x)$ 出发，考虑一族邻近函数
  \[
    u_\epsilon(x)=u(x)+\epsilon\eta(x),
    \qquad \eta(a)=\eta(b)=0.
  \]
  \begin{itemize}
    \item $\eta(x)$ 指定允许改变的方向。
    \item 小参数 $\epsilon$ 衡量改变的幅度。
    \item 固定端点要求 $\eta$ 在端点为零。
  \end{itemize}
  一阶变分定义为
  \[
    \delta J[u;\eta]=
    \left.\frac{\mathrm{d}}{\mathrm{d}\epsilon}J[u+\epsilon\eta]
    \right|_{\epsilon=0}.
  \]
  驻定路径满足 $\delta J=0$，并且该条件对每个可容许 $\eta$ 都成立。
\end{frame}

\begin{frame}{Euler--Lagrange 方程的推导}
  对 $J[u]=\int_a^bL(x,u,u')\,\mathrm{d}x$，
  \[
    \delta J=\int_a^b\left(L_u\eta+L_{u'}\eta'\right)\mathrm{d}x.
  \]
  对第二项分部积分：
  \[
    \delta J=\left[L_{u'}\eta\right]_a^b+
    \int_a^b\left(L_u-\frac{\mathrm{d}}{\mathrm{d}x}L_{u'}\right)\eta\,\mathrm{d}x.
  \]
  固定端点使边界项为零；$\eta$ 在内部任意，因此
  \[
    \boxed{L_u-\frac{\mathrm{d}}{\mathrm{d}x}L_{u'}=0.}
  \]
  \begin{block}{关键技巧}
    分部积分把导数从任意扰动 $\eta'$ 移到已知系数上，使任意性能够逐点发挥作用。
  \end{block}
\end{frame}

\begin{frame}{最短曲线的 Euler--Lagrange 解}
  对弧长泛函，
  \[
    L=\sqrt{1+u'^2},\qquad L_u=0,
    \qquad L_{u'}=\frac{u'}{\sqrt{1+u'^2}}.
  \]
  Euler--Lagrange 方程给出
  \[
    \frac{\mathrm{d}}{\mathrm{d}x}
    \left(\frac{u'}{\sqrt{1+u'^2}}\right)=0.
  \]
  因此 $u'$ 为常数，进而
  \[
    u''=0,\qquad u(x)=mx+b.
  \]
  \begin{alertblock}{结论}
    两点之间的直线最短，不只是几何直觉，也可视为弧长泛函驻定条件的结果。
  \end{alertblock}
\end{frame}

\begin{frame}{约束、力学与场论}
  \begin{columns}[T]
    \column{0.48\textwidth}
    \begin{block}{积分约束}
      若 $\int_a^bg(x,u,u')\,\mathrm{d}x=C$，引入拉格朗日乘子 $\lambda$，对
      \[
        L+\lambda g
      \]
      应用 Euler--Lagrange 方程。
    \end{block}
    \column{0.48\textwidth}
    \begin{block}{Hamilton 原理}
      真实轨道使作用量
      \[
        S=\int_{t_1}^{t_2}L(q,\dot q,t)\,\mathrm{d}t
      \]
      驻定，导出 Euler--Lagrange 运动方程。
    \end{block}
  \end{columns}
  对场 $u(\mathbf{x})$，变分方程推广为
  \[
    \frac{\partial F}{\partial u}-
    \partial_i\left(\frac{\partial F}{\partial(\partial_i u)}\right)=0.
  \]
  这正是连续介质和经典场论的共同语言。
\end{frame}
